%==============================================================================
% CAPÍTULO 18 — IA EM PERIODONTIA
% PARTE IV: APLICAÇÕES CLÍNICAS DA IA
%==============================================================================
\chapter{Inteligência Artificial em Periodontia}
\label{cap:periodontia-ia}
\niveltres

\begin{quote}
\textit{``A periodontite é a sexta doença mais prevalente no mundo, afetando 11,2\% da população global em sua forma grave. A inteligência artificial oferece, pela primeira vez, a possibilidade de triagem populacional com sensibilidade comparável à de especialistas.''}
\end{quote}

\section{Introdução: O Desafio Diagnóstico em Periodontia}

A doença periodontal representa um dos maiores desafios diagnósticos da odontologia contemporânea. Diferentemente da cárie dentária --- cuja detecção é predominantemente morfológica e localizada --- a periodontite envolve um complexo de fatores inflamatórios, microbiológicos e sistêmicos que demandam avaliação multidimensional. O periodontista precisa integrar profundidade de sondagem (PS), nível de inserção clínica (NIC), sangramento à sondagem (SS), perda óssea radiográfica, fatores de risco modificáveis (tabagismo, diabetes) e histórico de progressão para estadiar e graduar corretamente a doença.

A classificação de 2018 da \textit{American Academy of Periodontology} (AAP) e da \textit{European Federation of Periodontology} (EFP) introduziu um sistema de estadiamento (I--IV, baseado na severidade e complexidade) e graduação (A--C, baseado na taxa de progressão e fatores de risco) que, embora mais preciso, aumentou substancialmente a complexidade diagnóstica. Estudos mostram que a concordância interexaminador para esta classificação é moderada (kappa 0,45--0,65), mesmo entre especialistas. É neste contexto que a inteligência artificial emerge como ferramenta de suporte diagnóstico.

\citacaoativa{doi-1111-jcpe-13890}{10.1111/jcpe.13890}{A acurácia diagnóstica agrupada dos modelos de deep learning para detecção de perda óssea periodontal em radiografias panorâmicas atingiu sensibilidade de 87,5\% e especificidade de 90,2\%, com AUC média de 0,93, demonstrando desempenho comparável a periodontistas experientes.}

Neste capítulo, exploramos as aplicações da IA em periodontia em quatro eixos: (1) classificação automatizada segundo os critérios AAP/EFP 2018; (2) detecção e quantificação de perda óssea alveolar por deep learning; (3) modelos tabulares para predição de risco e prognóstico; e (4) integração com condições sistêmicas. Cada eixo é fundamentado em evidências publicadas e complementado por exemplos práticos utilizando a biblioteca \texttt{odontoia}.

\section{Classificação de Periodontite: Estágio e Grau com CNNs}

\subsection{A Classificação AAP/EFP 2018 e Seus Desafios}

A classificação de 2018 representa uma mudança paradigmática: de um sistema essencialmente descritivo (leve/moderada/grave) para um sistema multidimensional que captura a biologia da doença. O estadiamento (I--IV) reflete a severidade no momento do diagnóstico e a complexidade do tratamento, enquanto a graduação (A--C) incorpora a taxa de progressão histórica, fatores de risco e impacto sistêmico.

A complexidade operacional desta classificação é significativa. Para um único paciente, o clínico precisa:
\begin{itemize}
    \item Medir PS e NIC em seis sítios por dente (até 168 medições em uma boca completa);
    \item Avaliar perda óssea radiográfica (percentual de raiz, padrão horizontal ou vertical);
    \item Documentar envolvimento de furca (graus I--III);
    \item Calcular a razão perda óssea/idade;
    \item Integrar fatores de risco (tabagismo, HbA1c);
    \item Determinar se houve perda dentária por periodontite.
\end{itemize}

Não surpreende que a variabilidade interexaminador seja um problema documentado. A IA oferece padronização.

\subsection{CNNs para Estadiamento Automatizado}

O estudo de Park et al. (2023) representa um marco na classificação automatizada de periodontite. Utilizando uma arquitetura \textit{ResNet-50} modificada com \textit{attention gates} treinada em 4.000 dentes individuais com exames periodontais completos, o modelo classificou corretamente o estágio em 85,7\% dos casos (classificação em quatro classes: saúde, estágio I/II, estágio III, estágio IV).

\citacaoativa{doi-1016-j-jebdp-2024-102058}{10.1016/j.jebdp.2024.102058}{A CNN alcançou AUC de 0,96 para discriminação binária entre periodontite e saúde periodontal, com acurácia de 85,7\% para classificação em quatro estágios. As maiores confusões diagnósticas ocorreram entre estágios adjacentes (I vs. II e III vs. IV), fenômeno também observado em periodontistas humanos.}

\destaque{
\textbf{Interpretabilidade com Grad-CAM:} Os mapas de ativação revelaram que a CNN focalizou consistentemente três regiões-chave: (1) a crista óssea alveolar, especialmente em áreas com padrão de perda angular; (2) a região de furca em molares; e (3) o espaço do ligamento periodontal. Este comportamento é notável porque essas são precisamente as regiões que um periodontista treinado avalia prioritariamente, sugerindo que a rede aprendeu características clinicamente relevantes, não artefatos de imagem.
}

\subsection{Modelos Multimodais: Imagem + Dados Clínicos}

Uma limitação das abordagens puramente imagiológicas é que a periodontite é, por definição, uma doença inflamatória modulada por fatores sistêmicos. O trabalho de Costa et al. (2024) aborda esta limitação com um modelo \textit{gradient boosting} (XGBoost) que integra \textit{features} extraídas de radiografias panorâmicas por CNN com dados demográficos e clínicos tabulares (idade, sexo, tabagismo, diabetes mellitus).

\citacaoativa{doi-1007-s00784-024-05678-3}{10.1007/s00784-024-05678-3}{A inclusão de dados demográficos tabulares aumentou a AUC do modelo em 5\% comparado ao uso exclusivo de imagens radiográficas, alcançando AUC final de 0,92 para predição de periodontite. O tabagismo e o diabetes foram os fatores de risco mais impactantes no modelo integrado, com valores SHAP consistentes com a literatura epidemiológica.}

\subsubsection{Vantagens da Abordagem Multimodal}\indice{modelo multimodal}\indice{SHAP}\indice{triagem populacional}

O modelo multimodal oferece três vantagens sobre abordagens puramente radiográficas:
\begin{enumerate}
    \item \textbf{Maior acurácia}: a combinação de fontes de dados captura dimensões complementares da doença;
    \item \textbf{Explicabilidade}: os valores SHAP (\textit{SHapley Additive exPlanations}) quantificam a contribuição de cada fator de risco, permitindo auditoria clínica;
    \item \textbf{Aplicabilidade em triagem populacional}: o modelo não requer exame periodontal completo --- funciona com radiografia de rotina + questionário de saúde.
\end{enumerate}

\section{Detecção de Perda Óssea em Radiografias Panorâmicas}

\subsection{O Problema da Quantificação Objetiva}

A avaliação radiográfica da perda óssea periodontal é notoriamente subjetiva. Estudos clássicos de Benn~\cite{benn1992} (1992) demonstraram que a variabilidade intraexaminador na medição linear da distância da junção cemento-esmalte à crista óssea pode exceder 0,5 mm, valor clinicamente significativo considerando que a progressão anual média da periodontite é de 0,2--0,3 mm. A IA oferece quantificação objetiva e reprodutível.

\subsection{U-Net e YOLO-v9 para Perda Óssea}

A arquitetura U-Net, originalmente desenvolvida para segmentação de estruturas biomédicas, domina as aplicações de detecção de perda óssea periodontal. A revisão de Lee et al. (2024), abrangendo 19 estudos, confirma que modelos baseados em U-Net e Mask R-CNN consistentemente superam métodos tradicionais de processamento de imagem (limiarização, watershed) tanto em acurácia quanto em reprodutibilidade.

O trabalho de Ameli~\cite{ameli2024} (2024) representa o estado da arte na automação da medição de perda óssea. A abordagem combina U-Net para segmentação semântica do osso alveolar (identificando a crista óssea com precisão de pixel) com YOLO-v9 para detecção e contagem de dentes.

\subsubsection{Pipeline Automatizado para Medição Óssea}\indice{U-Net}\indice{YOLO-v9}\indice{pipeline}

O pipeline realiza:
\begin{enumerate}
    \item Detecção de cada dente individual com YOLO-v9 (\textit{bounding box} + classe dentária);
    \item Segmentação fina da junção cemento-esmalte e crista óssea com U-Net;
    \item Cálculo automático da distância JCE--crista óssea em milímetros (calibrado por escala radiográfica);
    \item Classificação do percentual de perda óssea por dente e categorização em estágios (I--IV).
\end{enumerate}

\destaque{
\textbf{Métrica de Performance~\cite{ameli2024} (Ameli 2024):} Dice score médio de 0,91 para segmentação de crista óssea; correlação de Pearson $r = 0,94$ entre medições automáticas e manuais por dois periodontistas independentes; tempo médio de processamento de 8,3 segundos por radiografia panorâmica (vs. 12--18 minutos para medição manual completa).
}

\subsection{Detecção de Envolvimento de Furca em 3D}

Uma limitação fundamental das radiografias 2D (panorâmica, periapical) é a incapacidade de avaliar adequadamente lesões de furca --- um determinante crítico de prognóstico e complexidade de tratamento. O estudo de Shin et al. (2023) utilizou uma rede \textit{3D CNN} (\textit{DenseVoxelNet}) para detecção automática de envolvimento de furca em CBCT, alcançando sensibilidade de 91,5\% e especificidade de 93,8\%.

\citacaoativa{doi-1038-s41598-019-44839-3}{10.1038/s41598-019-44839-3}{A rede 3D CNN detectou envolvimentos de furca grau I (incipientes) com acurácia de 82,3\%, superando significativamente a avaliação radiográfica 2D por periodontistas (63,5\%), e demonstrando que a terceira dimensão, combinada com IA, pode revelar lesões subclínicas antes que se tornem radiograficamente evidentes.}

\section{Modelos Tabulares para Predição de Risco Periodontal}

\subsection{Além da Radiografia: Dados Clínicos e Sistêmicos}

Enquanto as CNNs excelentes em tarefas de visão computacional, muitos determinantes do risco periodontal são de natureza tabular: idade, tabagismo (maços-ano), controle glicêmico (HbA1c), frequência de higiene oral, histórico familiar, medicações e condições sistêmicas. Modelos de \textit{machine learning} tabular --- \textit{random forest}, \textit{XGBoost}, \textit{LightGBM} e redes neurais \textit{feedforward} --- são particularmente adequados para estes dados.

A revisão abrangente de Khan et al. (2024) em \textit{Periodontology 2000} mapeou o panorama completo das aplicações de deep learning em periodontia, incluindo modelos preditivos de progressão:

\citacaoativa{doi-1111-prd-12567}{10.1111/prd.12567}{Modelos de machine learning tabular demonstraram capacidade de estratificar pacientes em três categorias de risco de progressão periodontal (baixo, moderado, alto) com AUC de 0,88--0,94, superando significativamente os escores de risco tradicionais como o Periodontal Risk Assessment (PRA) de Lang \& Tonetti, que tipicamente alcança AUC de 0,72--0,78.}

\subsection{Features Preditivas: O Que os Modelos Aprendem}

A análise de importância de \textit{features} (via SHAP ou permutação) consistentemente identifica os seguintes preditores principais de progressão de periodontite em modelos tabulares:

\begin{table}[H]
\centering
\caption{Preditores Principais de Progressão Periodontal em Modelos de ML}
\label{tab:preditores-periodontite}
\begin{tabular}{p{4cm}p{2cm}p{6cm}}
\toprule
\textbf{Preditor} & \textbf{Importância Relativa} & \textbf{Interpretação SHAP} \\
\midrule
Tabagismo atual (maços-ano) & 28\% & Maior impacto para $>10$ maços-ano; efeito dose-dependente \\
HbA1c (\%) & 22\% & Risco aumenta exponencialmente acima de 7,0\% \\
Profundidade de sondagem basal média & 18\% & PS média $>4$ mm é ponto de inflexão \\
Perda óssea radiográfica inicial ($>30\%$) & 15\% & Maior preditor de perda dentária futura \\
Sangramento à sondagem (\% sítios) & 10\% & SS $>25\%$ dos sítios indica atividade \\
Idade & 7\% & Interage com perda óssea na razão perda/idade \\
\bottomrule
\end{tabular}
\end{table}

\subsection{Treinando um Classificador de Risco Periodontal com \texttt{odontoia}}

A biblioteca \texttt{odontoia} oferece uma interface simplificada para treinamento de modelos de risco periodontal. O módulo \texttt{odontoia.models.dental\_ml} encapsula o pipeline completo de pré-processamento, treinamento e avaliação.

\subsubsection{Implementação do Classificador}\indice{classificador de risco}\indice{XGBoost}\indice{SHAP}

\begin{pratica}{Treinamento de Classificador de Risco Periodontal}
O exemplo abaixo utiliza dados simulados representativos de um estudo clínico transversal com 2.000 pacientes, incluindo profundidade de sondagem média, percentual de sítios com sangramento, índice de placa, tabagismo e HbA1c.

\begin{lstlisting}[language=Python, caption={Treinamento de classificador de risco periodontal com XGBoost}]
from odontoia.models.dental_ml import (
    PeriodontalRiskClassifier,
    PeriodontalDataset
)
from sklearn.model_selection import train_test_split
import numpy as np

# 1. Carregar dataset periodontal
dataset = PeriodontalDataset.from_csv("dados_periodontais.csv")
X = dataset.features   # PS, NIC, SS, placa, tabagismo, HbA1c
y = dataset.labels     # 0=baixo, 1=moderado, 2=alto risco

# 2. Split estratificado
X_train, X_test, y_train, y_test = train_test_split(
    X, y, test_size=0.2, stratify=y, random_state=42
)

# 3. Treinar classificador XGBoost
clf = PeriodontalRiskClassifier(
    model_type="xgboost",
    n_estimators=200,
    max_depth=6,
    learning_rate=0.05,
    use_shap=True
)
clf.fit(X_train, y_train)

# 4. Avaliar
accuracy = clf.score(X_test, y_test)
print(f"Acuracia: {accuracy:.3f}")

# 5. Explicabilidade com SHAP
shap_values = clf.explain(X_test.iloc[:10])
clf.plot_shap_summary(shap_values, feature_names=X.columns)
\end{lstlisting}

\textbf{Resultado esperado:}
\begin{lstlisting}[style=saida]
Acuracia: 0.892
AUC (one-vs-rest): [0.94, 0.89, 0.91]
Top-3 features SHAP: ['tabagismo_macos_ano', 'hbA1c', 'PS_media']
\end{lstlisting}
\end{pratica}

\teste{O classificador deve atingir acurácia $\ge 0,85$ em dados de teste, com AUC mínima de 0,88 para a classe de alto risco. O relatório SHAP deve identificar tabagismo e HbA1c entre os top-3 preditores.}

\section{Periodontite e Condições Sistêmicas: IA como Ponte}

\subsection{O Eixo Periodontal-Sistêmico}

A relação bidirecional entre periodontite e doenças sistêmicas --- particularmente diabetes mellitus tipo 2 e doenças cardiovasculares --- está bem estabelecida na literatura. A periodontite aumenta a resistência insulínica via citocinas pró-inflamatórias (TNF-$\alpha$, IL-6), enquanto a hiperglicemia crônica exacerba a inflamação periodontal via AGEs (\textit{Advanced Glycation End-products}) e estresse oxidativo.

A IA oferece ferramentas únicas para investigar estas interações complexas:
\begin{itemize}
    \item \textbf{Modelos de mediação}: quantificam vias causais entre variáveis (ex.: diabetes $\rightarrow$ inflamação $\rightarrow$ perda óssea);
    \item \textbf{Modelos de interação}: detectam efeitos sinérgicos entre fatores de risco (ex.: tabagismo $\times$ diabetes);
    \item \textbf{Estratificação de risco integrada}: combinam biomarcadores séricos, microbioma oral e parâmetros clínicos periodontais.
\end{itemize}

\subsection{Diabetes e Periodontite: Modelos de Predição Bidirecional}

O estudo de Natarajan et al.~\cite{natarajan2024} (2024) desenvolveu um modelo \textit{ensemble} (stacking de random forest + XGBoost + regressão logística regularizada) que prediz simultaneamente o risco de:
\begin{enumerate}
    \item Progressão de periodontite em pacientes diabéticos (AUC 0,91);
    \item Descontrole glicêmico (HbA1c $> 7,0\%$) em pacientes com periodontite (AUC 0,87);
    \item Necessidade de intervenção periodontal cirúrgica em 24 meses (AUC 0,89).
\end{enumerate}

O modelo utiliza 47 \textit{features} incluindo PS média, NIC, SS, número de dentes perdidos por periodontite, HbA1c atual e histórico, índice de massa corporal, tabagismo e PCR ultrassensível.

\subsection{Tabagismo: O Maior Fator de Risco Modificável}

Em todos os modelos de predição de risco periodontal, o tabagismo consistentemente emerge como o preditor mais impactante --- frequentemente com importância SHAP 2--3$\times$ maior que o segundo preditor. Modelos de IA são particularmente úteis para:
\begin{itemize}
    \item Quantificar o efeito dose-resposta (maços-ano) na perda óssea;
    \item Predizer a resposta ao tratamento em fumantes vs. não-fumantes;
    \item Identificar subgrupos de fumantes com perda óssea acelerada (fenótipo de ``high responder'').
\end{itemize}

\destaque{
\textbf{Implicação clínica:} Modelos de IA identificam que fumantes com $>20$ maços-ano e periodontite estágio III têm probabilidade 4,2$\times$ maior de perda dentária em 5 anos comparados a não-fumantes com mesma severidade, mesmo sob tratamento periodontal regular. Esta estratificação permite intensificar a frequência de manutenção periodontal.
}

\section{Estudos de Caso: IA Periodontal na Prática Clínica}

\subsection{Caso 1: Triagem Populacional com Modelo Multimodal}

Em um estudo multicêntrico brasileiro conduzido em 2024--2025, pesquisadores da Universidade de São Paulo (USP) e da Universidade Estadual de Campinas (UNICAMP) implantaram um sistema de triagem periodontal baseado em IA em 12 Unidades Básicas de Saúde (UBS) no estado de São Paulo. O sistema utilizava o modelo multimodal de Costa et al. (2024) --- que integra radiografia panorâmica (features extraídas por EfficientNet-B4) com dados demográficos (idade, sexo, tabagismo, diabetes) via XGBoost --- para classificar 4.500 pacientes em três categorias de risco periodontal (baixo, moderado, alto).

Os resultados preliminares demonstraram que:
\begin{itemize}
    \item O sistema identificou corretamente 92\% dos pacientes com periodontite estágio III--IV (alto risco), que foram encaminhados para atendimento especializado;
    \item A triagem automatizada reduziu em 67\% o tempo de espera para consulta com periodontista (de 4,2 meses para 1,4 meses em média);
    \item A concordância entre o modelo e a avaliação presencial de periodontistas (kappa 0,81) foi considerada substancial, validando o uso da ferramenta como filtro inicial.
\end{itemize}

\subsubsection{Impacto na Saúde Pública}\indice{triagem populacional}\indice{saúde pública}\indice{UBS}

Este caso ilustra o potencial da IA para odontologia de saúde pública, onde a demanda por especialistas em periodontia excede em muito a oferta, e a triagem automatizada pode otimizar o encaminhamento de casos prioritários.

\subsection{Caso 2: Monitoramento Longitudinal de Periodontite com Séries Radiográficas}

O estudo de Khan et al. (2024) relatou um sistema de monitoramento longitudinal que utiliza uma arquitetura \textit{siamese network} (rede siamesa) para comparar radiografias panorâmicas do mesmo paciente em dois momentos temporais (ex.: baseline e 12 meses). A rede aprende a detectar alterações sutis na crista óssea alveolar que podem indicar progressão ativa da doença antes que a perda óssea atinja limiares clinicamente evidentes.

Em uma coorte de 350 pacientes com periodontite estágio II--III sob manutenção periodontal trimestral, o sistema identificou 28 pacientes (8\%) com sinais precoces de progressão que não haviam sido detectados pelo periodontista na avaliação visual convencional. Destes, 22 (79\%) confirmaram progressão na consulta subsequente (com sondagem completa), demonstrando a capacidade da IA de detectar alterações subclínicas que escapam à inspeção visual humana.

\subsection{Caso 3: Assistência à Decisão Cirúrgica em Defeitos de Furca}

\notaexplicativa{A decisão de tratar um envolvimento de furca por meios não-cirúrgicos (raspagem e alisamento radicular) vs. cirúrgicos (acesso cirúrgico, regeneração, tunelização ou exodontia) é notoriamente complexa. Um estudo de Shin et al. (2023) demonstrou que o uso de IA para análise de CBCT pré-operatório alterou a decisão cirúrgica em 31\% dos casos (45/145 defeitos de furca), tipicamente indicando intervenção mais precoce e menos conservadora quando comparado ao planejamento baseado apenas em radiografias 2D.}

\section{Discussão: O Valor Incremental da IA em Periodontia}

A questão central que emerge da literatura não é mais \textit{se} a IA pode auxiliar o diagnóstico periodontal, mas \textit{qual o valor incremental} que ela agrega sobre a prática clínica padrão. A metanálise de Lee et al. (2024) estabeleceu que modelos de deep learning atingem sensibilidade de 87,5\% e especificidade de 90,2\% para detecção de perda óssea --- comparável a periodontistas experientes. Mas o valor real da IA não está em substituir o especialista, e sim em três funções complementares:

\begin{enumerate}
    \item \textbf{Triagem e priorização}: Em contextos de alta demanda (saúde pública, clínicas-escola, convênios), a IA pode processar centenas de radiografias em minutos, sinalizando casos de alto risco para avaliação prioritária pelo periodontista;

    \item \textbf{Segundo leitor}: Como ``segunda opinião automatizada'', a IA reduz erros de omissão --- lesões que o clínico deixa de notar por fadiga, distração ou viés de confirmação. Estudos mostram que a combinação dentista + IA consistentemente supera o dentista sozinho (GANho médio de sensibilidade de 8--12\%);

    \item \textbf{Quantificação objetiva}: A medição da perda óssea em milímetros, automatizada e reprodutível, elimina a subjetividade inerente à estimativa visual, permitindo monitoramento quantitativo da progressão ao longo do tempo --- essencial para decisões de \textit{timing} de intervenção.
\end{enumerate}

A evidência atual também sugere que modelos multimodais (imagem + dados tabulares) consistentemente superam modelos unimodais em periodontia, com ganho de AUC de 5--8\%. Isto é esperado, pois a periodontite é uma doença inflamatória modulada sistemicamente --- qualquer modelo puramente imagiológico é cego a fatores de risco críticos como tabagismo, controle glicêmico e predisposição genética.

\citacaoativa{doi-1111-prd-12567}{10.1111/prd.12567}{A incorporação de dados de microbioma oral (sequenciamento 16S rRNA) e biomarcadores salivares (MMP-8, IL-1$\beta$) como features adicionais em modelos de ML representa a próxima fronteira da periodontia de precisão. Modelos com features microbiológicas e inflamatórias atingiram AUC de 0,94--0,96 para predição de progressão em 24 meses, superando abordagens exclusivamente radiográficas (AUC 0,85--0,88).}

\section{Limitações Atuais e Lacunas de Evidência}

\subsection{Validade Externa e Generalização}

A quase totalidade dos estudos de IA em periodontia publicados até 2025 foram conduzidos em datasets de instituições únicas, predominantemente na Ásia (Coreia do Sul, China, Japão) e Europa. A validade externa destes modelos para populações latino-americanas, africanas e de diferentes estratos socioeconômicos permanece não testada. Diferenças na prevalência de periodontite agressiva, padrões de perda óssea, acesso a tratamento e comorbidades podem degradar significativamente a performance dos modelos.

\notaexplicativa{O fenômeno de \textit{dataset shift} --- quando a distribuição dos dados na população-alvo difere da distribuição nos dados de treinamento --- é particularmente preocupante em periodontia, onde a prevalência e severidade da doença variam dramaticamente entre populações. A periodontite grave (estágio III--IV) afeta 11,2\% da população global, mas esta proporção pode variar de 3\% (em coortes escandinavas com excelente acesso odontológico) a 25\% (em populações de baixa renda sem acesso a cuidados periodontais).}

\subsection{Falta de Padrão-Ouro Independente}

A maioria dos estudos utiliza a avaliação de periodontistas como padrão-ouro (\textit{ground truth}) para treinar modelos de IA. Entretanto, como extensamente documentado, a concordância interexaminador para classificação AAP/EFP 2018 é apenas moderada (kappa 0,45--0,65). Isto cria um paradoxo: estamos treinando IA para replicar um padrão-ouro que é intrinsecamente ruidoso. Abordagens que utilizam consenso de múltiplos especialistas ou, idealmente, desfechos clínicos objetivos (perda dentária em 5--10 anos) como \textit{ground truth} permanecem raras.

\subsection{Equidade Algorítmica em Periodontia}

\citacaoativa{doi-1007-s00784-024-05678-3}{10.1007/s00784-024-05678-3}{Modelos treinados predominantemente em populações asiáticas podem apresentar viés significativo quando aplicados a outras etnias. Fatores como dimensões craniofaciais médias, prevalência de biotipo gengival fino vs. espesso, e padrões de perda óssea (horizontal vs. vertical) variam entre grupos étnicos, potencialmente introduzindo disparidades de performance que precisam ser ativamente monitoradas e mitigadas.}

\section{Arquiteturas Comparadas para Periodontia}

\subsection{CNN 2D vs. CNN 3D vs. Transformer Híbrido}

A escolha da arquitetura de deep learning para tarefas periodontais depende da modalidade de imagem, do desfecho clínico-alvo e dos recursos computacionais disponíveis. A Tabela~\ref{tab:arquiteturas-periodontia} sintetiza as evidências comparativas.

\begin{table}[H]
\centering
\caption{Comparação de Arquiteturas de Deep Learning para Tarefas Periodontais}
\label{tab:arquiteturas-periodontia}
\begin{tabular}{p{2.5cm}p{2.2cm}p{2.2cm}p{2.5cm}p{2.5cm}}
\toprule
\textbf{Tarefa} & \textbf{U-Net (2D)} & \textbf{CNN 3D (CBCT)} & \textbf{Hybrid CNN-Transformer} & \textbf{Melhor Escolha} \\
\midrule
Perda óssea (panorâmica) & Dice 0,90--0,93 & N/A & Dice 0,91--0,94 & U-Net (melhor custo-benefício) \\
Furca (CBCT 3D) & Dice 0,65--0,72 & Dice 0,87--0,91 & Dice 0,85--0,89 & CNN 3D (necessita 3ª dimensão) \\
Classificação estágios & AUC 0,94--0,96 & AUC 0,93--0,95 & AUC 0,95--0,97 & Transformer (captura contexto global) \\
Segmentação crista óssea & Dice 0,89--0,92 & Dice 0,91--0,94 & Dice 0,92--0,95 & Híbrido (precisão máxima) \\
\bottomrule
\end{tabular}
\end{table}

\subsection{Pipeline Multimodal com Fusão Tardia}

O código abaixo implementa um pipeline completo de fusão tardia (\textit{late fusion}) para classificação de estágios de periodontite, combinando uma CNN (EfficientNet-B4) para extração de features radiográficas com um MLP para dados tabulares, seguindo a metodologia de Costa et al. (2024):

\subsubsection{Código do Pipeline}\indice{pipeline multimodal}\indice{fusão tardia}\indice{EfficientNet}

\begin{lstlisting}[language=Python, caption={Pipeline multimodal para classificação de periodontite com fusão tardia}]
import torch
import torch.nn as nn
from odontoia.models.dental_ml import (
    MultimodalPeriodontalClassifier,
    PeriodontalDataset
)
from odontoia.explainability import SHAPAnalyzer

# 1. Dataset multimodal
dataset = PeriodontalDataset.from_multimodal_csv(
    "periodontia_multimodal.csv",
    image_col="panoramica_path",
    tabular_cols=["idade", "sexo", "tabagismo_macos_ano", 
                   "diabetes", "hbA1c", "PS_media", "SS_pct"],
    label_col="estadio_periodontite"  # I, II, III, IV
)

# 2. Classificador com fusao tardia
clf = MultimodalPeriodontalClassifier(
    image_encoder="efficientnet-b4",  # Features radiograficas
    tabular_encoder="mlp",            # MLP para dados clinicos
    fusion_method="late_concat",      # Concatenacao tardia
    n_classes=4,
    use_attention_fusion=True         # Atencao cross-modal
)

# 3. Treinamento com metricas periodontais
clf.fit(
    dataset,
    epochs=50,
    class_weights="balanced",  # Compensa desbalanceamento
    metrics=["accuracy", "kappa", "f1_macro"]
)

# 4. Explicabilidade multimodal
analyzer = SHAPAnalyzer(clf)
shap_results = analyzer.explain_multimodal(
    dataset,
    n_samples=100,
    feature_names=dataset.feature_names
)

# 5. Relatorio de interpretabilidade
print("=== IMPORTANCIA SHAP MULTIMODAL ===")
print(f"Features de imagem: {shap_results.image_importance:.2%}")
print(f"Features tabulares: {shap_results.tabular_importance:.2%}")
print(f"\nTop-3 features tabulares:")
for feat, val in shap_results.top_tabular(3):
    print(f"  {feat}: SHAP = {val:.3f}")

# 6. Validacao cruzada estratificada por estagio
cv_results = clf.cross_validate(
    dataset, cv=5, stratify_by="estadio",
    verbose=True
)
print(f"\nCV Mean Accuracy: {cv_results['accuracy_mean']:.3f} "
      f"(+/- {cv_results['accuracy_std']:.3f})")
print(f"CV Mean Kappa:   {cv_results['kappa_mean']:.3f} "
      f"(+/- {cv_results['kappa_std']:.3f})")
\end{lstlisting}

\textbf{Resultado esperado:}
\begin{lstlisting}[style=saida]
=== IMPORTANCIA SHAP MULTIMODAL ===
Features de imagem: 64.3%
Features tabulares: 35.7%

Top-3 features tabulares:
  tabagismo_macos_ano: SHAP = 0.412
  hbA1c: SHAP = 0.298
  PS_media: SHAP = 0.187

CV Mean Accuracy: 0.903 (+/- 0.021)
CV Mean Kappa:   0.868 (+/- 0.031)
\end{lstlisting}

\teste{O classificador multimodal deve atingir acurácia $\ge 0,88$ e kappa $\ge 0,83$ na validação cruzada. A fusão tardia com atenção cross-modal deve mostrar melhoria de pelo menos 5\% de acurácia sobre o modelo unimodal (apenas imagem). As features tabulares (tabagismo, HbA1c) devem figurar entre os top-3 preditores SHAP.}

\section{Desafios e Limitações da IA em Periodontia}

\subsection{Qualidade e Padronização de Dados}

A periodontia enfrenta desafios específicos de dados:
\begin{itemize}
    \item \textbf{Variabilidade de sondagem}: A profundidade de sondagem é intrinsecamente variável (força de sondagem, angulação, tipo de sonda). Modelos treinados com dados de um examinador podem não generalizar para outros;
    \item \textbf{Datasets pequenos}: A maioria dos estudos publicados utiliza datasets com $< 500$ pacientes, limitando a generalização;
    \item \textbf{Desbalanceamento de classes}: Periodontite severa (estágios III--IV) representa $< 20\%$ da população geral, exigindo técnicas de \textit{oversampling} ou \textit{class weighting}.
\end{itemize}

\subsection{Gap de Validação Prospectiva}

Apesar dos resultados promissores em estudos retrospectivos, a validação prospectiva --- o verdadeiro teste de utilidade clínica --- permanece rara. Menos de 15\% dos estudos de IA em periodontia incluíram qualquer forma de validação externa ou prospectiva. Esta lacuna é crítica porque modelos podem aprender \textit{shortcuts} --- correlações espúrias específicas do dataset que não se sustentam em novos contextos clínicos.

\subsection{Integração com o Fluxo Clínico}

Para que a IA periodontal tenha impacto clínico real, precisa integrar-se ao fluxo de trabalho do consultório:

\subsubsection{Fluxo de Trabalho com IA Periodontal}\indice{fluxo clínico}\indice{integração}

\begin{enumerate}
    \item O paciente realiza radiografia panorâmica (já é rotina);
    \item A IA processa a imagem em segundos, gerando relatório de perda óssea por dente;
    \item O periodontista realiza sondagem confirmatória nos sítios sinalizados;
    \item O modelo multimodal calcula risco integrado e sugere intervalo de manutenção;
    \item O relatório é anexado ao prontuário eletrônico.
\end{enumerate}

\section{O Futuro: Periodontia de Precisão com IA}

A convergência de IA, sequenciamento de microbioma e dispositivos \textit{point-of-care} aponta para uma periodontia de precisão:
\begin{itemize}
    \item \textbf{Microbioma oral + ML}: Classificadores baseados em perfil microbiano (16S rRNA) predizem sítios com risco de progressão antes que a perda de inserção seja detectável;
    \item \textbf{Monitoramento contínuo}: Escovas inteligentes com sensores de pressão e câmeras intraorais acopladas a modelos de IA para detecção precoce de inflamação gengival;
    \item \textbf{Medicina periodontal personalizada}: Algoritmos que recomendam antibióticos específicos baseados no perfil microbiano e predição de resistência;
    \item \textbf{Gêmeos digitais periodontais}: Integração de dados longitudinais para simular trajetórias de progressão sob diferentes cenários de intervenção.
\end{itemize}

% ===================================================================
% SEÇÃO DE REFINAMENTO: Limitações Metodológicas
% Adicionar ao final dos capítulos técnicos (5, 7, 8, 11, 18)
% Auditoria multi-engine 2026
% ===================================================================

\section{Limitações Metodológicas e Ceticismo Estruturado}

\notaexplicativa{Esta seção é resultado direto da auditoria multi-engine
(11 motores) executada em 2026. Identifica premissas implícitas que
devem ser explicitadas para o leitor.}

% ------- PREMISSAS IMPLÍCITAS -------
\subsection*{Premissas Implícitas que Este Capítulo Assume}

\begin{enumerate}
    \item \textbf{Validade externa}: os resultados de acurácia
    citados são de estudos em populações específicas
    (China, Irã, EUA, Arábia). A \emph{miscigenação populacional brasileira}
    pode alterar a acurácia em 5-15\%.
    \item \textbf{Definição de verdade de campo}: a ``verdade''
    (ground truth) usada para treinar os modelos vem de anotação humana
    especialista, que tem \emph{variabilidade inter-observador}
    reportada em 8-15\% para cárie (Kappa de Cohen típico: 0.6-0.8).
    \item \textbf{Tipo de imagem}: resultados são específicos para
    radiografias panorâmicas e bitewing. Para periapical, CBCT ou
    foto intraoral, a acurácia tipicamente cai.
    \item \textbf{Threshold operacional}: o threshold de decisão
    padrão (0.5) não é o ótimo clínico. Cada centro deve recalibrar
    com base na prevalência local e no custo de falsos negativos vs
    positivos.
    \item \textbf{Data leakage}: alguns estudos na literatura
    podem ter vazamento entre treino e teste --- a meta-análise
    de Rohban (2024) \cite{rohban2024} identifica este risco em
    até 23\% dos estudos analisados.
\end{enumerate}

% ------- CRITÉRIOS DE BRADFORD-HILL -------
\subsection*{Avaliação de Causalidade (Critérios de Bradford-Hill)}

\notaexplicativa{Para qualquer afirmação do tipo ``X melhora Y'',
o motor causal do \texttt{multi\_reasoning} aplica os 9 critérios
de Bradford-Hill (1965).}

Para o claim central deste capítulo --- ``deep learning detecta
lesões com acurácia superior a clínicos'' --- os critérios
satisfatórios são:

\begin{itemize}
    \item[$\checkmark$] \textbf{Força da associação}: AUC médio 0.93 (forte).
    \item[$\checkmark$] \textbf{Consistência}: replicado em 87 estudos.
    \item[$\checkmark$] \textbf{Temporalidade}: avaliação sempre pós-treinamento.
    \item[$\sim$] \textbf{Plausibilidade}: modelo detecta padrões
    visuais; biologicamente plausível, mas não comprovado mecanismo.
    \item[$\times$] \textbf{Especificidade}: deep learning detecta
    \emph{muitas} coisas, não apenas cáries.
    \item[$\times$] \textbf{Gradiente}: não há dose-resposta estudada.
    \item[$\times$] \textbf{Experimento}: raros ensaios clínicos
    randomizados controlados.
    \item[$\times$] \textbf{Coesão}: específico para imagem odontológica?
    Incerto.
    \item[$\sim$] \textbf{Analogia}: paralelo com radiologia geral
    (Esteva 2017, Nature).
\end{itemize}

\notaexplicativa{Confiança atual: 5-6 de 9 critérios satisfeitos
$\approx 0.55$--$0.67$. O motor causal atribui confiança 0.7 (não
0.2) porque os critérios fortes (força, consistência, analogia)
estão presentes.}

% ------- RECOMENDAÇÕES PRÁTICAS -------
\subsection*{Recomendações para Implementação Clínica}

\begin{enumerate}
    \item \textbf{Validação local}: cada centro deve testar o modelo
    em pelo menos 200 casos locais antes do uso clínico.
    \item \textbf{Calibração}: usar \texttt{odontoia.metrics.classification.youden\_index}
    para encontrar o threshold ótimo.
    \item \textbf{Monitoramento}: reportar acurácia continuamente
    (drift monitoring). Acurácia decai quando o equipamento de
    imagem muda.
    \item \textbf{Fallback}: sempre ter revisão humana como
    segunda opinião. IA é apoio, não substituta.
    \item \textbf{Documentação}: manter o \emph{Model Card} atualizado
    conforme o framework FUTURE-AI \cite{futureai2025}.
\end{enumerate}

% ------- REFERÊNCIA À AUDITORIA -------
\notaexplicativa{Esta seção segue o protocolo recomendado pela
auditoria multi-engine do OpenCode Ecosystem. Para mais detalhes,
ver o Apêndice E: Ceticismo Estruturado.}


\section{Síntese do Capítulo}

A inteligência artificial está transformando a periodontia em quatro frentes complementares:

\begin{enumerate}
    \item \textbf{Classificação automatizada}: CNNs alcançam acurácia de 85--87\% na classificação de estágios da periodontite (AAP/EFP 2018), com AUC de 0,96 para discriminação saúde/doença. Os mapas de ativação (Grad-CAM) confirmam que as redes aprendem características clinicamente relevantes --- crista óssea, furca, espaço do ligamento periodontal.
    
    \item \textbf{Detecção de perda óssea}: Modelos baseados em U-Net e YOLO-v9~\cite{ameli2024} (Ameli 2024) quantificam perda óssea com correlação $r = 0,94$ com medições manuais, reduzindo o tempo de análise de 15 minutos para 8 segundos. A extensão para CBCT 3D permite detecção de lesões de furca incipientes com acurácia 82,3\%.
    
    \item \textbf{Predição de risco multimodal}: Modelos que integram dados radiográficos e clínicos (AUC 0,92) superam escores de risco tradicionais (AUC $\sim$0,75). O tabagismo e o diabetes emergem consistentemente como os preditores mais impactantes.
    
    \item \textbf{Ponte sistêmica}: A IA quantifica interações complexas entre periodontite e condições sistêmicas, permitindo estratificação de risco personalizada e intervenções preventivas direcionadas.
\end{enumerate}

As limitações atuais --- datasets pequenos, validação prospectiva insuficiente, integração clínica incipiente --- são desafios transitórios. A trajetória é clara: a periodontia caminha para um modelo de precisão onde algoritmos de IA, integrados ao fluxo clínico, auxiliam o periodontista em cada etapa: do diagnóstico ao prognóstico, da decisão terapêutica ao monitoramento longitudinal.

\subsection*{Leituras Recomendadas}
\begin{itemize}
    \item Khan et al. (2024). Deep learning-based periodontal disease assessment: a comprehensive review. \textit{Periodontology 2000}. \url{https://doi.org/10.1111/prd.12567}
    \item Lee et al. (2024). Deep learning for detection of periodontal bone loss in panoramic radiographs. \textit{Journal of Clinical Periodontology}. \url{https://doi.org/10.1111/jcpe.13890}
    \item Costa et al. (2024). Machine learning models for prediction of periodontitis using panoramic radiographs and demographic data. \textit{Clinical Oral Investigations}. \url{https://doi.org/10.1007/s00784-024-05678-3}
    \item Park et al. (2023). Convolutional neural network for classification of periodontitis stages. \textit{Journal of Periodontology}. \url{https://doi.org/10.1016/j.jebdp.2024.102058}
\end{itemize}


% ===================================================================
% SEÇÃO DE REFINAMENTO: Contexto Brasileiro (SUS + SB Brasil)
% Adicionar ao Cap 18 (Periodontia) e Cap 25 (Ética)
% ===================================================================

\section{Contexto Brasileiro: SB Brasil, SUS e Populações Locais}

\notaexplicativa{A auditoria multi-engine identificou que vários
claims de impacto (``500K profissionais'', ``democratização'')
careciam de ancoragem no sistema de saúde brasileiro real.}

\subsection*{Dados Epidemiológicos Nacionais (SB Brasil 2020)}

A Pesquisa Nacional de Saúde Bucal (SB Brasil 2020) \cite{sbbrasil2020}
é o maior levantamento epidemiológico odontológico do Brasil. Os
dados mais relevantes para este capítulo (periodontia) são:

\begin{table}[H]
\centering
\small
\begin{tabularx}{\textwidth}{lXcc}
\toprule
\textbf{Faixa Etária} & \textbf{Condição} & \textbf{Prevalência} & \textbf{Comparação 2010} \\
\midrule
12 anos & Sangramento gengival & 35,4\% & +8\% \\
15-19 anos & Cálculo dentário & 28,2\% & +5\% \\
35-44 anos & Perda de inserção $\geq$ 4mm & 31,3\% & +12\% \\
65-74 anos & Edentulismo (perda total) & 43,2\% & -3\% \\
\bottomrule
\end{tabularx}
\caption{Prevalência de condições periodontais no Brasil (SB Brasil 2020).}
\label{tab:sbbrasil-periodontia}
\end{table}

\subsection*{Implicações para Modelos de IA}

\notaexplicativa{Modelos treinados em populações asiáticas ou
europeias podem subestimar a prevalência e severidade de doença
periodontal no Brasil.}

\begin{enumerate}
    \item \textbf{Calibração}: modelos devem ser recalibrados para a
    população brasileira usando dados locais. O pacote
    \texttt{odontoia.metrics.classification.youden\_index} facilita isso.
    \item \textbf{Validação externa}: a métrica de um modelo treinado
    em dados iranianos (p.ex., Sharififar 2024) pode cair 10-15\%
    quando aplicado a pacientes do SUS de Maceió-AL.
    \item \textbf{Justiça algorítmica}: verificar se o modelo não
    performa pior em pacientes com menor acesso a tratamento
    periodontal prévio (proxy: região geográfica).
\end{enumerate}

\subsection*{O Caso SISTEGES: IA + Periodontia no SUS}

\notaexplicativa{Desde 2014, o SISTEGES (Sistema de
Telemonitoramento de Gestantes em Odontologia) opera no SUS de
Maceió-AL. É um caso real de IA aplicada à periodontia em
saúde pública brasileira.}

O SISTEGES integra:
\begin{itemize}
    \item Questionário de risco periodontal aplicado via app.
    \item Triagem automática baseada em regras fuzzy (motor
    \texttt{fuzzy} do \texttt{multi\_reasoning}).
    \item Encaminhamento para UBS com classificação de risco
    (baixo/médio/alto).
    \item Acompanhamento longitudinal de marcadores inflamatórios.
\end{itemize}

Desde sua implementação, o SISTEGES cobriu mais de 5.000 gestantes
e demonstrou redução de 23\% na incidência de periodontite
gestacional não tratada (dados internos 2024).

\subsection*{Recomendações para Implementação no SUS}

\begin{enumerate}
    \item \textbf{Integração}: usar o TOTVS-GOV (sistema do DATASUS)
    como backend, evitando fragmentação.
    \item \textbf{Capacitação}: 1 dentista por UBS como
    ``embaixador IA'' --- mais efetivo que treinamento massivo.
    \item \textbf{Avaliação}: incluir o protocolo do Cap.~15
    (Grad-CAM) para verificar explicabilidade em decisões clínicas.
    \item \textbf{Custo}: estimar R\$ 50.000-100.000/ano por município
    médio (10-15 UBS).
\end{enumerate}
